\section{Introduction}

I/O memory protection prevents errant direct memory access (DMA)
    from buggy or malicious devices and is essential
    to safety and security of virtualized environments.
Modern virtual machines (VMs) 
    primarily use passthrough devices 
%    through direct device assignment, 
%    Single Root I/O Virtualization (SR-IOV)~\cite{ms_sriov_overview,intel_sriov_ethernet,amd_sriov_mxgpu},
%    and Scalable I/O Virtualization (Scalable IOV)~\cite{intel_scalable_iov_spec}
    for near-native I/O performance,
    and thus desire both inter- and intra-VM I/O memory protection.
Inter-VM protection ensures a device (or a virtual function) to only access memory 
    regions explicitly authorized for the assigned VM;
    it prevents unauthorized DMA % from buggy or malicious devices
    into memory of other VMs. 
The protection is achieved by using an Input/Output Memory Management Unit (IOMMU)
    to isolate host I/O memory.

Intra-VM I/O memory protection, on the other hand, 
    enables fine-grained DMA protections for devices (or virtual functions) assigned to the VM,
    preventing unrestricted DMA 
    to the entire physical address space of the VM~\cite{Amit2011vIOMMU,RajapakshaIOMMU,AlexCodeInjection}.
Intra-VM protection is commonly supported by virtualized IOMMUs (vIOMMUs)~\cite{Amit2011vIOMMU,Tian2020coIOMMU}. 
A vIOMMU exposes an emulated IOMMU to the guest kernel;
the guest can directly use vIOMMUs with no kernel modification or hardare 
    support.
\tianyin{put a note on Jiyuan's vfio point---can we benefit vfio?}

% These mechanisms allow virtual machines (VMs) to interact with hardware devices 
%    with minimal hypervisor intervention, 
%    significantly improving I/O throughput and latency. 
    
As memory protection is on the critical path of I/O operations,
    it has significant impacts on application performance.
Prior work has extensively optimized IOMMU-enable data paths by 
    minimizing IOTLB misses~\cite{?} and 
    reducing I/O page-table walk overhead~\cite{Rubin2024FastSafe,markuze2018damn}.
%    extensive efforts have been made 
%    to minimize IOMMU overhead, making
%    I/O memory protection fast and safe~\cite{Rubin2024FastSafe,markuze2018damn}.
While these efforts eliminate IOMMU overhead 
%However, we find that existing efforts, which target host IOMMU,
    for inter-VM protection in virtualized environments;
    the overhead of intra-VM protection using vIOMMU remains substantial.
% , intra-VM I/O memory protection with vIOMMU could 
%    incur significant performance overheads~\cite{Amit2011vIOMMU,Tian2020coIOMMU,Lv2022LAvIOMMU}.
Our measurement \tianyin{provide necessary details; now it's too hard to 
    interpret this number} 
    reveals that enabling strict I/O protection with vIOMMU \tianyin{details}
    introduces up to 90\% throughput degradation.
Such an overhead is unaffordable to most I/O-intensive applications
    and forces reduced protection, \textit{e.g.}, 
    Linux offers the \texttt{lazy} mode of I/O memory protection
    which defer IOTLB invalidation to optimize performance,
    at the cost of reduced device isolation (\S\ref{sec:security}).

In this paper, we show that
    vIOMMU-based intra-VM memory protection manifests fundamentally different
    performance bottleneck (compared with 
    IOMMU for inter-VM protection or bare-mental environments),
    and desire new principles of performance optimizations.
% the performance of I/O memory protection
%    without compromising its safety and security.
Specifically, with optimized address
    translation~\cite{Rubin2024FastSafe,markuze2018damn},
    the bottleneck of vIOMMU-enabled protection is dominated by the overhead of
    IOTLB invalidation which is required whenever the device driver unmaps a memory page
    and must be done before the page is freed and reused.
Compared with IOMMU-based protection on bare mental,
    IOTLB invalidation in virtualization environments has to additionally
    triggers VM exists, because the physical IOMMU invalidation queue is a shared host resource
    and is managed by the host kernel.
Compared with hardware invalidation (which
    takes 200--300 nanoseconds on modern hardware~\cite{Rubin2024FastSafe}),
    VM exits are orders of magnitude more costly (10--20 microseconds).
Unfortunately, existing vIOMMU system designs~\cite{Amit2011vIOMMU,Tian2020coIOMMU,linux} 
    do not effectively mitigate the cost of invalidation-induced VM exists, 
    but negatively amplify its effect into a compounding overhead.

Concretely, to provide strong safety guarantee, existing systems 
    (\textit{e.g.}, \texttt{strict} mode
    in Linux) make the invalidation 
    a {\it synchronous} process: the OS writes an invalidation request 
    to the invalidation queue 
    as soon as the device driver unmaps a memory page 
    and blocks the running thread until the IOMMU completes the invalidation.
% IOTLB invalidation must be done before the unmapped page is freed and reused
    % as soon as the vIOMMU unmaps a DMA-ed page to 
 %   to prevent the device from using a stale translation and breaking memory protection.
Moreover, the OS strictly {\it serializes} concurrent 
    unmap operations from different cores; inflight unmap operations 
    are blocked if there is an ongoing invalidation.
    % blocking future invalidation requests before the ongoing invalidatoin finishes,
    % This serialization caps throughput independent of the core count.
As a result, the IOMMU overhead grows super-linearly with the number of cores 
    due to lock contention \tianyin{which lock?}.
Optimizations that use asynchronous or deferred invalidation are developed, but all with 
    reduced safety~\cite{Amit2011vIOMMU}, \textit{e.g.}, deferred invalidation
    opens a window that exposes stale translations,
    breaking memory protection.

% Intuitively, a single VM exit can consume thousands of CPU cycles,
% whereas a native memory operation takes only tens.
% Consequently, the per-unmap request latency increases drastically in a VM.
% Yet, despite this massive latency spike per request,
% the system's \emph{parallelism} remains artificially capped by legacy software synchronization,
% preventing the system from hiding this latency.
% Our profiling of the recent Linux kernel~\cite{linux_v6129}---which already adopts DAMN's page pool optimization---
% shows that even with hardware-assisted nested IOMMU,
% turning on vIOMMU still suffers up to 82\% throughput degradation
% on the Tx datapath compared with no vIOMMU.

% Recent architectures have developed dedicated support for IOMMU virtualization to address these performance challenges.
% Early vIOMMU implementations utilized shadow page tables, which suffered from high overhead due to the frequent VM exits required for IOTLB invalidation on every mapping and unmapping operation.
% These VM exits were necessary to notify the hypervisor and host kernel of each change\cite{Amit2011vIOMMU,qemu_vtd}.
% To improve performance, recent architectures introduced hardware-assisted Nested IOMMU support, 
%     which eliminates VM exits when the guest creates a new IOMMU mapping.
% However, this approach still requires IOTLB invalidation---and thus a VM exit---at unmap time~\cite{yiliu1765qemu}.
% This is because the IOMMU hardware is shared by all VMs, 
% and allowing multiple VMs to independently write to the invalidation queue would lead to corruption~\cite{intel_vtd_spec_caching, qemu_intel_iommu_c}.

We argue that existing designs of vIOMMU-enabled memory protection
%    \tianyin{[be precise about what are existing designs...in Linux? drivers?]}
    unnecessarily conflate \emph{memory-safety ordering} 
    with \emph{cross-core serialization}.
Strict protection (the strongest safety property) only requires
    correct ordering
    for a \emph{specific} memory mapping within a single thread:
    for a mapping from IO Virtual Address (IOVA) to guest physical address (GPA),
    the IOVA must be unmapped, then invalidated in the hardware IOTLB, 
    and only then the physical page can be freed and reused.
But, it does \emph{not} require serializing independent unmap operations 
    across different cores.
With this insight, we redesign vIOMMU-based 
    I/O memory protection mechanisms for virtualized environments
    to \tianyin{what's our performance goal?}
    while ensuring the strongest safety of strict protection.
Our key principle is to enable the system to process highly concurrent 
    IOTLB invalidation requests and thus amortize VM exit costs.
We realize this principle in a system named \brand 
    through the following endevours.

First, we present \emph{\optOne}, a technique that decouples
    the submission of IOTLB invalidation requests from the actual invalidation process.
\tianyin{it's hard to catch the essence of the idea}
A delegated core performs the IOTLB invalidation work,
while other cores submit invalidation requests and wait.
This decoupling removes the need to wait for previous invalidation 
    to finish before generating a new invalidation request,
    removing mitigating the cross-core serialization bottleneck.
Furthermore, it can process concurrent invalidation requests through efficient batching:
multiple invalidation requests can be coalesced
while the delegated core is processing a previous batch,
    effectively reducing the total number of VM exits.
\emph{\optOne} requires zero changes to device drivers;
    all modifications are contained within the IOMMU subsystem
    in the OS kernel.
\tianyin{need an argument why this is safe}

\tianyin{Where is the distributed handling to eliminate the single 
    core bottleneck?}

% To address these challenges, we introduce \brand that 
% significantly improve vIOMMU performance while maintaining security parity with the strict model.
% First, we developed \emph{\optOne}, a technique that decouples the submission of invalidation requests from the actual invalidation process.
% A delegated core will perform the IOTLB invalidation work, while other cores submit invalidation requests
% and wait for their requests to finish.
% The decoupling of request submission and invalidation also reduces lock contention, 
%     as generating an invalidation request no longer requires holding the global lock.
% What's more, the design enables efficient batching of invalidation requests, 
%     as multiple invalidation requests can be coalesced together while the delegated core is processing a previous batch.
% The coalescence of invalidation requests reduce the number of VM exits, and thus significantly improves performance.
% \emph{\optOne} can be realized with zero change to any drive code; 
%     all changes are contained within the IOMMU subsystem, making it broadly applicable to all drivers.
%  offloads the invalidation process to a delegated core, 

% By decoupling the request submission from the invalidation process, we ensure that only the enqueue operation remains under . 
% By moving the IOTLB invalidation out of the locking protection, we significantly reduce lock contention, leaving only the enqueue operation under lock. 
% This design also enables \emph{Efficient Batching}, as multiple invalidation requests can be grouped together while the delegated core is processing a previous batch.

% However, to maintain the strict security model without sacrificing this newfound parallelism,
% the driver must not free the associated memory page
% until the asynchronous invalidation completes;
% otherwise, a Use-After-Free vulnerability arises~\cite{RajapakshaIOMMU}.

Second, 
we present the \emph{Deferred Free Page (DFP)} mechanism
    to enforce strict safety ordering, % (unmap $\rightarrow$ invalidate $\rightarrow$ free),
    without forcing the CPU to synchronously wait,
% To maintain the strict security model, 
% the driver must wait for the invalidation to complete \emph{before} 
%    freeing the associated guest physical address (GPA) page.
% Otherwise, a the guest OS is vulnerable to data confidentiality attack\cite{RajapakshaIOMMU}.
DFP introduces a new callback-based interface in the IOMMU subsystem,
    allowing drivers to pass a callback function
    that delays physical memory cleanup until hardware invalidation is confirmed.
With DFP, the driver can proceed immediately without waiting for the invalidation to complete
    while still preserving strict security guarantees.
% The cominbation of DFP and \optOne can enable the system to process more DMA packets since it spends less time in IOMMU subsystem.
% and, thus, improve the throughput without sacrificing strict security guarantees.
The cominbation of DFP and \optOne{} can eliminate most of CPU blocking time
    in the IOMMU subsystem, achieving close-to maximum throughput.

% However, simply delegating invalidation is not enough. 
% To maintain the strict security model, 
% the driver must wait for the invalidation to complete \emph{before} freeing the associated memory page to prevent data integrity and confidentiality issues.
% Drivers assume that a page that is logically unmapped is also physically unmapped, 
%     so the drivers may release the physical page before the hardware invalidation actually finishes.
% This can lead to security issues if the page is reused for a different purpose (e.g. store password) before the invalidation completes,
%     as the device could still access the page through DMA~\cite{RajapakshaIOMMU}.
% We remove the waiting time without lose of strict garrantee through
%     \emph{Deferred Free Page (DFP)} mechanism.
% Refactoring all drivers to handle this is impractical due to missing IOMMU subsystem interface and engineering scale. 
% To solve this, we propose the \emph{Deferred Free Page (DFP)} mechanism. 
% DFP introduces a new callback-based interface in the IOMMU subsystem.
% This allows drivers to pass a callback function and delay the physical memory cleanup until the hardware invalidation is confirmed;
%     as a result, the driver can proceed without waiting for the invalidation to complete, while still maintaining the strict security model.
% This technique is broadly applicable and compatible with both lazy and \optOne modes.
% With DFP combined with \optOne, we can achieve close-to-no-vIOMMU performance while maintaining the strict security model.

We implement \brand on top of Linux kernel 6.12.9
with fewer than 1.5k lines of code changed (only 373 lines in the device driver).
\brand achieves close-to-no-vIOMMU performance
while maintaining the strict security model,
yielding up to a 7x throughput improvement over the baseline strict model.

In summary, we make the following contributions:
\begin{itemize}[leftmargin=*]
    \item A detailed analysis of the inefficiencies in existing vIOMMU designs,
    identifying cross-core serialization as the root cause of performance bottlenecks;
    \item \brand as a new design that achieves near-native performance
        by enabling the system to process highly concurrent invalidation requests and minimize blocking time in the IOMMU subsystem;
    \item  A security analysis of \brand which shows that \brand achieves the strongest 
        safety property that can be offered by IOMMU;   
%    non-strict policies
%        and \brand achieves the highest level of security without sacrificing performance.
%    \item \textbf{Evaluations on modern hardware:} We evaluate \brand on the latest hardware
%    with nested translation support using high-end 400 Gbps networking workloads.
    \item We release our research artifacts to support reproducibility.
\end{itemize}

\if 0
\para{vIOMMU}

The I/O memory management unit (IOMMU) provides hardware-based protection against DMA attacks 
    by adding a level of indirection for DMA addresses. 
The IOMMUs virtualizes the address space of devices by providing it with an I/O Virtual Address (IOVA) 
    which it then translates to the Physical Address (PA). 
This way, the device can access only the pages explicitly allowed by the OS. 
Similar to the Memory Management Unit (MMU), the IOMMU uses page tables 
    located in main memory for address translation 
    and an IOTLB to cache recent accesses. 
The page tables are managed by the OS and as with the MMU, 
    have a page granularity \cite{AlexCodeInjection}. 
In virtualized environments, an IOMMU properly managed by the hypervisor 
    provides inter-VM protection by ensuring a VM can only access memory it is assigned. 
% Another benefit of IOMMU is that it enables legacy 32-bit devices to access addresses beyond their DMA capabilities.

A virtual IOMMU (vIOMMU)\cite{Amit2011vIOMMU} presented to a guest OS extends this provide intra-VM protection 
    allowing a VM to protect itself against malicious or buggy drivers. 
The address translation is split into two stages: the guest  manages the first stage, 
    which maps Guest Virtual Address (GVA) to Guest Physical Address (GPA) 
    and the hypervisor manages the second stage (nested translation), 
    which maps GPA to Host Physical Address (HPA). 
The hypervisor can implement nested translation either in software using shadow page tables, 
    or in hardware using modern extensions like Intel VT-d or AMD-Vi. 
\fi 